% GENERATED FILE. Edit canonical lessons, then run npm run generate:books.
% canonical-source-hash: fnv1a64:fc1b9a2fe4de89d1
% canonical-lessons: PA-C347-dhol, PA-C347-vajauna, PA-C347-tal, PA-C347-thap, PA-C347-tari

\chapter{The Big Drum}
\label{ch:pa-a2-the-big-drum}
\hlchaptermodality{347}

\begin{chapteropening}
\textbf{By the end of this chapter:} \emph{I can say that I like the drum and its beat.}

Complete the last lesson of chapter 347: I can say that I like the drum and its beat.
\end{chapteropening}

\section[\texorpdfstring{\pa{ਢੋਲ} (ḍhol)}{ਢੋਲ (ḍhol)}]{\pa{ਢੋਲ} (ḍhol) — a drum; a big two-sided drum}

\label{lesson:PA-C347-dhol}



\begin{quote}
Before the new one: say the Punjabi for helpless, then the Punjabi for a pity.
\end{quote}

\subsection*{You'll want to know: \pa{ਢੋਲ}}
\textbf{\pa{ਢੋਲ}} — \emph{ḍhol} — \textquotedblleft{}a drum; a big two-sided drum\textquotedblright{}.

\begin{quote}
\pa{ਮੈਨੂੰ} \pa{ਢੋਲ} \pa{ਦੀ} \pa{ਆਵਾਜ਼} \pa{ਬਹੁਤ} \pa{ਪਸੰਦ} \pa{ਹੈ।}
\end{quote}

\begin{quote}
\emph{mainū̃ ḍhol dī āvāz bahut pasand hai.}
\end{quote}

\begin{quote}
I really like the sound of the drum.
\end{quote}

A masculine noun.

\begin{cousinweb}
\textbf{\pa{ਢੋਲ}} (\emph{ḍhol}) is inherited from Sanskrit \emph{ḍhola}, \textquotedblleft{}drum\textquotedblright{}; CDIAL lists Punjabi \emph{ḍhol}, \textquotedblleft{}drum\textquotedblright{}, under it (CDIAL 5608).
\end{cousinweb}

\begin{grammarlens}[title={What you've built}]
One more word for this chapter.
\end{grammarlens}

\subsection*{Your turn}
\begin{itemize}
\raggedright
  \item \emph{Say it:} \emph{ḍhol}
  \item \emph{Say it:} \emph{ḍhol}, once more
  \item \emph{Say it:} \emph{mainū̃ ḍhol dī āvāz bahut pasand hai.} — I really like the sound of the drum.
  \item \emph{Recall:} say the Punjabi for helpless, then the Punjabi for a pity, then say \emph{ḍhol} again
\end{itemize}

\begin{tcolorbox}[breakable,colback=teal!4,colframe=teal!35!black,title={Before you move on}]
What does \pa{ਢੋਲ} mean? (\textquotedblleft{}A drum; a big two-sided drum\textquotedblright{}.) Say it once more.
\end{tcolorbox}
\section[\texorpdfstring{\pa{ਵਜਾਉਣਾ} (vajāuṇā)}{ਵਜਾਉਣਾ (vajāuṇā)}]{\pa{ਵਜਾਉਣਾ} (vajāuṇā) — to play an instrument; to beat a drum}

\label{lesson:PA-C347-vajauna}



\begin{quote}
Before the new one: say the Punjabi for a pity, then the Punjabi for a drum.
\end{quote}

\subsection*{You'll want to know: \pa{ਵਜਾਉਣਾ}}
\textbf{\pa{ਵਜਾਉਣਾ}} — \emph{vajāuṇā} — \textquotedblleft{}to play an instrument; to beat a drum\textquotedblright{}.

\begin{quote}
\pa{ਮੇਰਾ} \pa{ਭਰਾ} \pa{ਸ਼ਾਮ} \pa{ਨੂੰ} \pa{ਢੋਲ} \pa{ਵਜਾਉਂਦਾ} \pa{ਹੈ।}
\end{quote}

\begin{quote}
\emph{merā bharā shām nū̃ ḍhol vajāundā hai.}
\end{quote}

\begin{quote}
My brother plays the drum in the evening.
\end{quote}

Past (I …): \pa{ਮੈਂ} \pa{ਢੋਲ} \pa{ਵਜਾਇਆ}, \emph{maiṁ ḍhol vajāiā}.

\begin{cousinweb}
\textbf{\pa{ਵਜਾਉਣਾ}} (\emph{vajāuṇā}) is the causative partner of \emph{vajjṇā}, \textquotedblleft{}to sound\textquotedblright{}, which goes back to Sanskrit \emph{vādyate}, \textquotedblleft{}is made to sound\textquotedblright{}; CDIAL lists Punjabi \emph{vajjṇā} under it (CDIAL 11513).
\end{cousinweb}

\begin{grammarlens}[title={What you've built}]
One more word for this chapter.
\end{grammarlens}

\subsection*{Your turn}
\begin{itemize}
\raggedright
  \item \emph{Say it:} \emph{vajāuṇā}
  \item \emph{Say it:} \emph{vajāuṇā}, once more
  \item \emph{Say it:} \emph{merā bharā shām nū̃ ḍhol vajāundā hai.} — My brother plays the drum in the evening.
  \item \emph{Recall:} say the Punjabi for a pity, then the Punjabi for a drum, then say \emph{vajāuṇā} again
  \item \emph{Write it:} copy \textbf{\pa{ਵਜਾਉਣਾ}} once, checking the whole word
\end{itemize}

\begin{tcolorbox}[breakable,colback=teal!4,colframe=teal!35!black,title={Before you move on}]
What does \pa{ਵਜਾਉਣਾ} mean? (\textquotedblleft{}To play an instrument; to beat a drum\textquotedblright{}.) Say it once more.
\end{tcolorbox}
\section[\texorpdfstring{\pa{ਤਾਲ} (tāl)}{ਤਾਲ (tāl)}]{\pa{ਤਾਲ} (tāl) — rhythm; the beat}

\label{lesson:PA-C347-tal}



\begin{quote}
Before the new one: say the Punjabi for a drum, then the Punjabi for to play an instrument.
\end{quote}

\subsection*{You'll want to know: \pa{ਤਾਲ}}
\textbf{\pa{ਤਾਲ}} — \emph{tāl} — \textquotedblleft{}rhythm; the beat\textquotedblright{}.

\begin{quote}
\pa{ਉਹ} \pa{ਤਾਲ} \pa{ਨਾਲ} \pa{ਨੱਚਦੀ} \pa{ਹੈ।}
\end{quote}

\begin{quote}
\emph{oh tāl nāl nachchdī hai.}
\end{quote}

\begin{quote}
She dances in time with the beat.
\end{quote}

A masculine noun.

\begin{cousinweb}
\textbf{\pa{ਤਾਲ}} (\emph{tāl}) is inherited from Sanskrit \emph{tāla}, \textquotedblleft{}clap, rhythm\textquotedblright{}; CDIAL lists Punjabi \emph{tāl}, \textquotedblleft{}beating time with the hands\textquotedblright{}, under it (CDIAL 5748).
\end{cousinweb}

\begin{grammarlens}[title={What you've built}]
One more word for this chapter.
\end{grammarlens}

\subsection*{Your turn}
\begin{itemize}
\raggedright
  \item \emph{Say it:} \emph{tāl}
  \item \emph{Say it:} \emph{tāl}, once more
  \item \emph{Say it:} \emph{oh tāl nāl nachchdī hai.} — She dances in time with the beat.
  \item \emph{Recall:} say the Punjabi for a drum, then the Punjabi for to play an instrument, then say \emph{tāl} again
  \item \emph{Write it:} copy \textbf{\pa{ਤਾਲ}} once, checking the whole word
\end{itemize}

\begin{tcolorbox}[breakable,colback=teal!4,colframe=teal!35!black,title={Before you move on}]
What does \pa{ਤਾਲ} mean? (\textquotedblleft{}Rhythm; the beat\textquotedblright{}.) Say it once more.
\end{tcolorbox}
\section[\texorpdfstring{\pa{ਥਾਪ} (thāp)}{ਥਾਪ (thāp)}]{\pa{ਥਾਪ} (thāp) — a drumbeat; a beat of the hand on a drum}

\label{lesson:PA-C347-thap}



\begin{quote}
Before the new one: say the Punjabi for to play an instrument, then the Punjabi for rhythm.
\end{quote}

\subsection*{You'll want to know: \pa{ਥਾਪ}}
\textbf{\pa{ਥਾਪ}} — \emph{thāp} — \textquotedblleft{}a drumbeat; a beat of the hand on a drum\textquotedblright{}.

\begin{quote}
\pa{ਢੋਲ} \pa{ਦੀ} \pa{ਥਾਪ} \pa{ਦੂਰ} \pa{ਤੱਕ} \pa{ਜਾਂਦੀ} \pa{ਹੈ।}
\end{quote}

\begin{quote}
\emph{ḍhol dī thāp dūr tak jāndī hai.}
\end{quote}

\begin{quote}
The beat of the drum carries a long way.
\end{quote}

A feminine noun.

\begin{cousinweb}
\textbf{\pa{ਥਾਪ}} (\emph{thāp}) is inherited from a reconstructed \emph{*thapp}, \textquotedblleft{}slap, pat\textquotedblright{}; CDIAL lists Punjabi \emph{thāp}, \textquotedblleft{}pat, slap\textquotedblright{}, under it and marks it masculine (CDIAL 6091). CDIAL does not give the sense \textquotedblleft{}drumbeat\textquotedblright{}, and the feminine use taught here (the beat of the drum) is not recorded there.
\end{cousinweb}

\begin{grammarlens}[title={What you've built}]
One more word for this chapter.
\end{grammarlens}

\subsection*{Your turn}
\begin{itemize}
\raggedright
  \item \emph{Say it:} \emph{thāp}
  \item \emph{Say it:} \emph{thāp}, once more
  \item \emph{Say it:} \emph{ḍhol dī thāp dūr tak jāndī hai.} — The beat of the drum carries a long way.
  \item \emph{Recall:} say the Punjabi for to play an instrument, then the Punjabi for rhythm, then say \emph{thāp} again
  \item \emph{Write it:} copy \textbf{\pa{ਥਾਪ}} once, checking the whole word
\end{itemize}

\begin{tcolorbox}[breakable,colback=teal!4,colframe=teal!35!black,title={Before you move on}]
What does \pa{ਥਾਪ} mean? (\textquotedblleft{}A drumbeat; a beat of the hand on a drum\textquotedblright{}.) Say it once more.
\end{tcolorbox}
\section[\texorpdfstring{\pa{ਤਾੜੀ} (tāṛī)}{ਤਾੜੀ (tāṛī)}]{\pa{ਤਾੜੀ} (tāṛī) — a clap; clapping}

\label{lesson:PA-C347-tari}



\begin{quote}
Before the new one: say the Punjabi for rhythm, then the Punjabi for a drumbeat.
\end{quote}

\subsection*{You'll want to know: \pa{ਤਾੜੀ}}
\textbf{\pa{ਤਾੜੀ}} — \emph{tāṛī} — \textquotedblleft{}a clap; clapping\textquotedblright{}.

\begin{quote}
\pa{ਗੀਤ} \pa{ਤੋਂ} \pa{ਬਾਅਦ} \pa{ਸਭ} \pa{ਲੋਕ} \pa{ਤਾੜੀਆਂ} \pa{ਵਜਾਉਂਦੇ} \pa{ਹਨ।}
\end{quote}

\begin{quote}
\emph{gīt ton bāad sabh lok tāṛīāṁ vajāunde han.}
\end{quote}

\begin{quote}
After the song everyone claps.
\end{quote}

A feminine noun.

\begin{cousinweb}
\textbf{\pa{ਤਾੜੀ}} (\emph{tāṛī}) is inherited from Sanskrit \emph{tāla}, \textquotedblleft{}clap, rhythm\textquotedblright{}; CDIAL lists Punjabi \emph{tālī}, \textquotedblleft{}clapping\textquotedblright{}, and Lahnda \emph{tāṛī}, \textquotedblleft{}clapping the hands\textquotedblright{}, under it (CDIAL 5748).
\end{cousinweb}

\begin{grammarlens}[title={What you've built}]
One more word for this chapter.
\end{grammarlens}

\subsection*{Your turn}
\begin{itemize}
\raggedright
  \item \emph{Say it:} \emph{tāṛī}
  \item \emph{Say it:} \emph{tāṛī}, once more
  \item \emph{Say it:} \emph{gīt ton bāad sabh lok tāṛīāṁ vajāunde han.} — After the song everyone claps.
  \item \emph{Recall:} say the Punjabi for rhythm, then the Punjabi for a drumbeat, then say \emph{tāṛī} again
  \item \emph{Write it:} copy \textbf{\pa{ਤਾੜੀ}} once, checking the whole word
\end{itemize}

\begin{tcolorbox}[breakable,colback=teal!4,colframe=teal!35!black,title={Before you move on}]
What does \pa{ਤਾੜੀ} mean? (\textquotedblleft{}A clap; clapping\textquotedblright{}.) Say it once more.
\end{tcolorbox}
